\documentclass{report}
\usepackage[letterpaper,portrait,margin=2cm]{geometry}
\usepackage{amsmath,amssymb,amsthm,mathtools,dsfont,xfrac,mathrsfs}
\usepackage{enumitem}
\usepackage[french]{babel}
\usepackage{titlesec}
\usepackage{fancyhdr}
\usepackage{hyperref}
\usepackage{tikz}
\usetikzlibrary{positioning, calc, automata}

\hypersetup{
	colorlinks=false
}

\allowdisplaybreaks

\fancyfoot{}
\fancyhead{}
\fancyhead[l]{MAT342 - Th\'eorie des anneaux}
\fancyhead[r]{\leftmark}
\fancyfoot[c]{\thepage}
\renewcommand{\headrulewidth}{.4pt}
\renewcommand{\footrulewidth}{0pt}

\setcounter{tocdepth}{3}

\pagestyle{fancy}

\title{MAT342 - Th\'eorie des anneaux

	Donn\'e par Juan Carlos Bustamante}
\author{Julien Houle}
\date{Automne 2026}

\renewcommand{\thesection}{ \arabic{chapter}.\arabic{section}}
\renewcommand{\thesubsection}{}
\renewcommand{\thesubsubsection}{}

\titleformat{\chapter}[hang]{\bfseries\huge\centering}{Chapitre \arabic{chapter}}{1em}{}[]
\titleformat{\section}[hang]{\bfseries\large}{Section\thesection}{1em}{}[]
\titleformat{\subsection}[hang]{\bfseries\normalsize}{}{0pt}{}[]
\titleformat{\subsubsection}[hang]{\slshape\normalsize}{}{0pt}{}[]

\newtheorem*{thm}{Th\'eor\`eme}
\newtheorem*{lem}{Lemme}
\newtheorem*{prop}{Proposition}
\newtheorem*{coro}{Corollaire}
\theoremstyle{definition}
\newtheorem*{dfn}{D\'efinition}
\theoremstyle{remark}
\newtheorem*{exem}{Exemple}
\newtheorem*{exer}{Exercice}
\newtheorem*{nota}{Notation}
\newtheorem*{rmq}{Remarque}
\newtheorem*{rapp}{Rappel}

\setlist{nosep}

\newenvironment*{subproof}[1][D\'emonstration]
{
\begin{proof}[#1]

}
{
\renewcommand{\qedsymbol}{$\blacksquare$}
\end{proof}
\renewcommand{\qedsymbol}{$\square$}
}

\def\lte{\leqslant}
\def\gte{\geqslant}

\def\NN{\mathbb{N}}
\def\ZZ{\mathbb{Z}}
\def\QQ{\mathbb{Q}}
\def\RR{\mathbb{R}}
\def\CC{\mathbb{C}}

\def\id{\mathds{1}}

\newcommand*{\abs}[1]{\left|#1\right|}
\newcommand*{\caract}[1]{\operatorname{car}\left(#1\right)}

\begin{document}
	\maketitle
	\pagenumbering{roman}
	\tableofcontents
	\newpage
	\pagenumbering{arabic}
	\chapter{Anneaux}
	\phantomsection\addcontentsline{toc}{section}{Rappel}
	\section*{Rappel}
	Un groupe est un ensemble avec une op\'eration $(G,\cdot)$ qui v\'erifient
	\begin{itemize}
		\item associativit\'e: $(xy)z=x(yz)$,
		\item $\exists$ neutre: $e_G$ t.q. $xe=ex=x$,
		\item Tout \'el\'ement est inversible: $\forall x, \exists x'$ t.q. $xx'=x'x=e$.
	\end{itemize}

	De plus, $G$ est ab\'elien si
	\begin{itemize}
		\item commutatif: $xy=yx$.
	\end{itemize}

	\begin{dfn}
		Un anneau est un ensemble muni de 2 op\'erations $(A,+,\cdot)$, t.q.
		\begin{enumerate}
			\item $(A,+)$ un groupe ab\'elien (neutre $0_A$),
			\item $(A,\cdot)$ un mono\"ide: assoc, neutre ($1_A$),
			\item Relation de compatibilit\'e
			\begin{itemize}
				\item Distributivit\'e: $x(y+z)=xy+xz$ et $(x+y)z=xz+yz$.
			\end{itemize}
		\end{enumerate}
	\end{dfn}

	\begin{rmq}
		\begin{enumerate}
			\item Certains textes ne n\'ecessite pas l'existence de neutre multiplicatif,
			\item $A$ est ``commutatif'' si l'op\'eration $\cdot$ l'est.
		\end{enumerate}
	\end{rmq}

	\begin{exem}
		\begin{enumerate}[label=\alph*)]
			\item $\ZZ$, $\QQ$, $\RR$, $\CC$ avec op\'erations usuelles;
			\item $\ZZ[x]$, $\QQ[x]$, $\RR[x]$, $\CC[x]$, $\ZZ[x,y]$, \dots: polyn\^omes;
			\item $M_n(\ZZ) = \ZZ^{n \times n}$: matrices carr\'ees $n \times n$ avec coefficients dans $\ZZ$, $M_n(\QQ)$, $M_n(\RR)$, $M_n(\CC)$, $M_n(\RR[x])$, \dots;
			\item $\ZZ_n = \sfrac{\ZZ}{n\ZZ}$ est un anneau commutatif,
			\item supposons $G$ un groupe ab\'elien: $\operatorname{end}(G) = \left\{ G \xrightarrow{f} G ~\middle|~ f\text{ homomorphisme de groupes} \right\}$, $f+g: G \to G, x \mapsto (f+g)(x) = f(x)+g(x)$, neutre $0$: la fonction constante \'egale \`a $0_G$, associativit\'e: h\'erit\'ee de $G$, ``oppos\'e'': $-f:x \mapsto -f(x) \in G$, alors $(\operatorname{end}(G), +)$ est un groupe ab\'elien. La composition est une autre op\'eration dans $\operatorname{end}(G)$,
			\begin{itemize}
				\item associativit\'e: oui,
				\item neutre: $\id_G: x \mapsto x$,
				\item distributivit\'e: pour $x \in G$ arbitraire,
				\begin{align*}
					(f\circ (g+h))(x)&= f((g+h)(x))\\
					&= f(g(x)+h(x))\\
					\shortintertext{comme $f$ est un homomorphisme de groupes}
					&= f(g(x))+f(h(x))\\
					&= (f \circ g)(x) + (f \circ h)(x)
				\end{align*}
				Ceci \'etant vrai pour tout $x \in G$, on a bien $f \circ (g+h) = f \circ g + f \circ h$ dans $\operatorname{end}(G)$. De m\^eme, on a la distributivit\'e de l'autre c\^ot\'e.
			\end{itemize}
			Ainsi, $\operatorname{end}(G)$ est un anneau, g\'en\'eralement non commutatif.
			\item Si $A,B$ sont des anneaux, $A \times B$ est un anneau avec $(a,b)+(a',b') \coloneq (a+a',b+b')$, $(a,b)\cdot(a',b') \coloneq (aa',bb')$, $\id_{A \times B} \coloneq (\id_A,\id_B)$;
			\item $X$ un ensemble, $A = \RR^{(X)}$: $\left\{ f:X \to \RR ~\middle|~ f\text{ est de support fini} \right\}$, c'est-\`a-dire $f(x)=0$ sauf pour un ensemble fini d'\'el\'ements.
			\begin{itemize}
				\item addition dans $A$: $f+g : X \to \RR : x \mapsto f(x)+g(x)$ fait de $(A,+)$ un groupe ab\'elien;
				\item multiplication dans $A$: $f \cdot g : X \to \RR : x \mapsto f(x) \cdot g(x)$ est associative et commutative. Les distributivit\'es tiennent. Le neutre multiplicatif serait la fonction constante $x \mapsto \id_\RR$ qui n'est pas \`a support fini.
			\end{itemize}
		\end{enumerate}
	\end{exem}

	\begin{prop}
		\begin{enumerate}
			\item Dans un anneau, le neutre multiplicatif est unique,
			\item Si un \'el\'ement poss\`ede un inverse multiplicatif, alors cet inverse est unique. L'inverse de $a$ est not\'e $a^{-1}$.
		\end{enumerate}
		\begin{proof}
			\begin{enumerate}
				\item Supposons $\id,\id'$ sont deux neutres multiplicatifs.
				\begin{align*}
					\id&= \id\id'\\
					\shortintertext{puisque $\id'$ est neutre}
					&= \id'\\
					\shortintertext{puisque $\id$ est neutre}
				\end{align*}
				\item Soit $a \in A$, $a',a''$ deux inverses.
				\begin{align*}
					a'&= a'\id\\
					\shortintertext{puisque $a''$ est inverse,}
					&= a'(aa'')\\
					\shortintertext{par associativit\'e,}
					&= (a'a)a''\\
					\shortintertext{puisque $a'$ est inverse,}
					&= \id a''\\
					&= a''
				\end{align*}
			\end{enumerate}
		\end{proof}
	\end{prop}

	\begin{prop}[R\`egles de calcul]
		Soit $(A,+,\cdot)$ un anneau.

		Pour tous $a,b,c \in A$, $m,n \in \ZZ$, on a
		\begin{enumerate}[label=\roman*)]
			\item $-(-a)=a$,
			\item $-(a+b)=-a-b$,
			\item $a+x=b+x \Rightarrow a=b$,
			\item \begin{align*}
				(m+n)a&= ma+na\\
				m(a+b)&= ma+mb\\
				m(na)&= (mn)a
			\end{align*}
			\item $a0=0a=0$,
			\item $a(-b)=(-a)b=-(ab)$,
			\item \begin{align*}
				a(b-c)&= ab-ac\\
				(a-b)c&= ac-bc
			\end{align*}
			\item $(-a)(-b)=ab$,
			\item \begin{align*}
				a\sum_{i=1}^{n} b_i&= \sum_{i=1}^{n} ab_i\\
				\left( \sum_{i=1}^{n} a_i \right)b&= \sum_{i=1}^{n} a_ib
			\end{align*}
			\item $(a^m)^n = a^{mn}$, $m,n \in \mathbb{N}$.
		\end{enumerate}

		\begin{subproof}[id\'ee de d\'emo]
			i) \`a iv): \c Ca se passe dans le groupe ab\'elien $(A,+)$, voir MAT141

			v): exercice

			vi): \begin{align*}
				a(-b)+ab&= a(-b+b)\\
				&= a0\\
				&= 0
			\end{align*}
			Il s'en suit que $a(-b)=-(ab)$. Idem pour l'autre \'egalit\'e.

			vii): exercice

			viii): \begin{align*}
				(-a)(-b)&= -((-a)b)\\
				&= -(-(ab))\\
				&= ab
			\end{align*}
			ix),x): r\'ecurrences.
		\end{subproof}
	\end{prop}

	\begin{dfn}
		$A$ un anneau donn\'e. Un sous-anneau de $A$ est un sous-ensemble $B$ de $A$ qui est un anneau pour les m\^emes op\'erations, et avec le m\^eme neutre multiplicatif: $\id_B=\id_A$. On \'ecrit $B \lte A$
	\end{dfn}

	\begin{exem}
		\begin{enumerate}
			\item $\QQ \lte \RR \lte \CC$.
			\item $A=\left\{ \begin{bsmallmatrix}
				a&0\\0&0
			\end{bsmallmatrix} ~\middle|~ a \in \RR \right\}$ est un anneau dont le neutre multiplicatif est $\begin{bsmallmatrix}
				1&0\\0&0
			\end{bsmallmatrix}$. Ce n'est pas un sous anneau de $\RR^{2 \times 2}$, car $\id_{\RR^{2\times2}} = \begin{bsmallmatrix}
				1&0\\0&1
			\end{bsmallmatrix} \notin A$.
			\item $\ZZ[i] = \left\{ a+bi ~\middle|~ a,b \in \ZZ \right\} \lte \CC$. Ce sont les entiers de Gauss, avec les m\^emes op\'erations que $\CC$.
			\item $d \in \ZZ$ pas divisible par un carr\'e. $\ZZ[\sqrt{d}] = \left\{ a+b\sqrt{d} ~\middle|~ a,b \in \ZZ \right\}$ avec les op\'erations induites par celles de $\RR$. $\ZZ[\sqrt{d}] \lte \QQ[\sqrt{d}] \lte \RR \lte \CC$.
			\item $T_n(\RR) \lte M_n(\RR)$, les matrices triangulaires inf\'erieures.
		\end{enumerate}
	\end{exem}

	\begin{prop}
		Soient $A$ un anneau et $B \subseteq A$. Alors, $B \lte A$ si, et seulement si,
		\begin{enumerate}
			\item $B$ est un sous-groupe ab\'elien de $A$,
			\item $a,b \in B \Rightarrow ab \in B$,
			\item $\id_A \in B$.
		\end{enumerate}
		\begin{proof}
			\begin{subproof}[$\Rightarrow$:]
				directe.
			\end{subproof}
			\begin{subproof}[$\Leftarrow$:]
				La multiplication est associative dans $A$, donc dans $B$ aussi. De m\^eme, on a les distributivit\'es. Le neutre $\id_A$ est neutre dans $A$, donc dans $B$ aussi.
			\end{subproof}
		\end{proof}
	\end{prop}

	\begin{exem}
		\begin{enumerate}
			\item Sous-anneaux (propres) de $\ZZ$: il n'y en a pas.

			$\ZZ_n$ n'est m\^eme pas un sous-ensemble de $\ZZ$.
			\item sous-anneaux de $\ZZ_m$: uniquement $\ZZ_m$.
			\item $A$ un anneau donn\'e, $Z(A)$ est $Z(A) = \left\{ a \in A ~\middle|~ ax=xa, \forall x \in A \right\}$. Voyons que $Z(A) \lte A$.
			\begin{itemize}
				\item $Z(A)$ est un sous-groupe additif de $A$. $0_A \in Z(A)$. De plus, si $a,a' \in Z(A)$ alors pour $x \in A$ on a
				\begin{align*}
					(a-a')x&= ax-a'x\\
					&= xa-xa'\\
					&= x(a-a')
				\end{align*}

				Donc, $Z(A)$ est un sous-groupe additif de $A$.
				\item $a,a' \in Z(A)$.
				\begin{align*}
					(aa')x&= a(a'x)\\
					&= a(xa')\\
					&= (ax)a'\\
					&= (xa)a'\\
					&= x(aa')
				\end{align*}

				Donc, $a,a' \in Z(A) \Rightarrow aa' \in Z(A)$.
				\item $\id_A \in Z(A)$, directement.
			\end{itemize}
			\begin{enumerate}
				\item $A$ commutatif $\Leftrightarrow Z(A)=A$.
				\item $A=\RR^{2\times2}$. Clairement, $\left\{ \lambda I = \begin{bsmallmatrix}
					\lambda&0\\0&\lambda
				\end{bsmallmatrix} ~\middle|~ \lambda\in\RR \right\} \in Z(A)$. Voyons que c'est tout. Si $\begin{bsmallmatrix}
					0&0\\\mu&0
				\end{bsmallmatrix} \in Z(A)$, alors
				\begin{align*}
					\begin{bsmallmatrix}
						0&0\\\mu&0
					\end{bsmallmatrix} \begin{bsmallmatrix}
						a&b\\c&d
					\end{bsmallmatrix} &= \begin{bsmallmatrix}
						a&b\\c&d
					\end{bsmallmatrix} \begin{bsmallmatrix}
						0&0\\\mu&0
					\end{bsmallmatrix}\\
					\begin{bsmallmatrix}
						0&0\\\mu a&\mu b
					\end{bsmallmatrix} &= \begin{bsmallmatrix}
						\mu b&0\\\mu d&0
					\end{bsmallmatrix}
				\end{align*}

				De m\^eme, aucune matrice $\begin{bsmallmatrix}
					0&\mu\\0&0
				\end{bsmallmatrix}$ avec $\mu\neq0$ ne se trouve dans $Z(A)$.

				Par ailleurs, $\begin{bsmallmatrix}
					\lambda&0\\0&\mu
				\end{bsmallmatrix} \notin Z(A) \Leftrightarrow \lambda\neq\mu$.
			\end{enumerate}
		\end{enumerate}
	\end{exem}

	\section{L'anneau de groupe}
	Soient $A$ un anneau commutatif et $G$ un groupe.

	\begin{rapp}
		$A^G = \left\{ f:G \to A \right\}$ toutes les fonctions est un anneau pour
		\begin{align*}
			(f_1+f_2)(x)&= f_1(x)+f_2(x)\\
			(f_1f_2)(x)&= f_1(x)f_2(x)
		\end{align*}
		avec les op\'erations induites par $A$.

		$A^{(G)} = \left\{ f:G \to A ~\middle|~ f\text{ est \`a support fini} \right\}$.
	\end{rapp}

	\begin{dfn}
		Pour $x \in G$, on pose $\delta_x : G \to A : y \mapsto \delta_x(y) = \left\{ \begin{array}{lll}
			\id_A&\text{si}&x=y\\
			0_A&\text{si}&x\neq y
		\end{array} \right.$. Si $f \in A^{(G)}$, alors $f(*) = \sum_{x \in G} f(x) \delta_x(*)$. En effet, pour $g_0 \in G$, $\sum_{x \in G} f(x) \delta_x(g_0) = f(g_0)\delta_{g_0}(g_0) = f(g_0)$.
	\end{dfn}

	\begin{rmq}
		$\left\{ \delta_x ~\middle|~ x \in G \right\}$ est une ``base'' pour l'ensemble $A^{(G)}$.
	\end{rmq}

	\begin{rmq}
		On a une injection $G \to A^{(G)} : x \mapsto \delta_x$ de sorte qu'on peut identifier $\delta_x$ de $A^{(G)}$ \`a $x$.
	\end{rmq}

	$A^{(G)}$ est un groupe ab\'elien pour l'addition induite par celle de $A$. Posons $\delta_x\bullet\delta_y \coloneq \delta_{xy}$ et on prolonge par distributivit\'e \`a $A^{(G)}$.
	\begin{align*}
		\left( \sum_{x \in G} a_xx \right) \left( \sum_{y \in G} b_yy \right)&= \sum_{x,y \in G} a_xb_y(xy)
	\end{align*}

	\begin{exem}
		\begin{enumerate}
			\item $A=\ZZ$, $G=\mathbb{D}_3 = \{\varepsilon, \rho, \sigma, \alpha, \beta, \gamma\}$, dont la table est pr\'esent\'ee ci-bas.
			\begin{table}[h]
				\centering
				$\begin{array}{c||c|c|c|c|c|c}
					\circ&\varepsilon&\alpha&\beta&\gamma&\rho&\sigma\\
					\hline\hline
					\varepsilon&\varepsilon&\alpha&\beta&\gamma&\rho&\sigma\\
					\hline
					\alpha&\alpha&\varepsilon&\rho&\sigma&\beta&\gamma\\
					\hline
					\beta&\beta&\sigma&\varepsilon&\rho&\gamma&\alpha\\
					\hline
					\gamma&\gamma&\rho&\sigma&\varepsilon&\alpha&\beta\\
					\hline
					\rho&\rho&\gamma&\alpha&\beta&\sigma&\varepsilon\\
					\hline
					\sigma&\sigma&\beta&\gamma&\alpha&\varepsilon&\rho
				\end{array}$
				\caption{Table de $\mathbb{D}_3$}
			\end{table}
			\begin{align*}
				\underbrace{\left( \rho+\sigma-\varepsilon \right)}_{\in \ZZ^{\mathbb{D}_3}} \underbrace{\left( 2\alpha-\beta \right)}_{\in \ZZ^{\mathbb{D}_3}}&= 2\rho\alpha + 2\sigma\alpha - 2\varepsilon\alpha - \rho\beta - \sigma\beta + \varepsilon\beta\\
				&= 2\gamma + 2\beta - 2\alpha - \alpha - \gamma + \beta\\
				&= \gamma + \beta - 3\alpha
			\end{align*}

			L'\'el\'ement neutre de $A^{(G)}$ est $\delta_e$, o\`u $e$ est le neutre de $G$.
			\item $A = \RR$, $G=\left\langle x \right\rangle = \left\langle x^n ~\middle|~ n \in \ZZ \right\rangle \cong \ZZ$. Les \'el\'ements de $A^{(G)}$ sont $\sum a_nx^n$, o\`u la somme est finie. Les polyn\^omes en $x$ et $x^{-1}$, par exemple $3x^{-7}+x^{-1}+4+\pi x^{25}$. Le produit: $x^nx^m=x^{n+m}$ donne lieu au produit traditionnel des polyn\^omes. En fait, $\RR[x] \lte A^{(G)}$.
			\item $A=\RR$, $G=\left\{ z \in \CC ~\middle|~ z^n=1 \right\} = \langle\omega\rangle = \{1,\omega,\omega^2,\cdots,\omega^{n-1}\}$, o\`u $\omega=e^{i\sfrac{2\pi}{n}}$. $\RR^{(G)}$ contient des \'el\'ements de la forme $a_0 + a_1\omega + a_2\omega^2 + \cdots a_{n-1}\omega^{n-1}$.
			\begin{align*}
				(1-\omega)(1+\omega+\omega^2+\cdots+\omega^{n-1})&= (1+\omega+\omega^2+\cdots+\omega^{n-1}) - (\omega+\omega^2+\cdots+\omega^n)\\
				&= 1-\omega^n\\
				&= 0
			\end{align*}
		\end{enumerate}
	\end{exem}

	\phantomsection\addcontentsline{toc}{section}{Alg\`ebre de chemins}
	\section*{Alg\`ebre de chemins}
	\begin{dfn}
		Un \emph{carquois} Q est un quadruplet $Q = (Q_0, Q_1, s, b)$, o\`u $Q_0,Q_1$ sont des ensembles, points et fl\`eches respectivement, et $s,b: Q_1 \to Q_0$ sont les applications ``source'' et ``but''.

		Si $\alpha\in Q_1$, on dessine $s(\alpha) \xrightarrow{\alpha} b(\alpha)$.
	\end{dfn}

	\begin{exem}
		\begin{figure}[h]
			\centering
			\label{figQ}
			\begin{tikzpicture}
				\node[state] (1) at (0,0) {1};
				\path[->] (1) edge[loop above] node {$\alpha$} (1);
				\node () [below=of 1] {$Q^1$};
			\end{tikzpicture}
			\hspace{15mm}
			\begin{tikzpicture}
				\node[state] (1) at (0,0) {1};
				\node[state] (2) [left=of 1] {2};
				\node[state] (3) [left=of 2] {3};
				\path[->]	(3) edge[bend left] node[above] {$\beta$} (2)
				(3) edge[bend right] node[below] {$\gamma$} (2)
				(2) edge node[above] {$\alpha$} (1);
				\node () [below=of 2] {$Q^2$};
			\end{tikzpicture}
			\hspace{15mm}
			\begin{tikzpicture}
				\node[state] (1) at (0,0) {1};
				\path[->] (1) edge[loop above] node[above] {$\alpha$} (1)
				(1) edge[loop below] node[below] {$\beta$} (1);
				\node () [below=of 1] {$Q^3$};
			\end{tikzpicture}
			\hspace{15mm}
			\begin{tikzpicture}
				\node[state] (1) at (0,0) {1};
				\node[state] (2) [right=of 1] {2};
				\node[state] (3) [right=of 2] {3};
				\path[->] (3) edge node[above] {$\beta$} (2)
				(2) edge node[above] {$\alpha$} (1);
				\node () [below=of 2] {$Q^4$};
			\end{tikzpicture}
			\caption{4 exemples de carquois}
		\end{figure}
	\end{exem}

	\begin{dfn}
		Un \emph{chemin} dans $Q$ est:
		\begin{itemize}
			\item un chemin stationnaire $\varepsilon_i$ pour chaque $i \in Q_0$,
			\item une suite de fl\`eches $\alpha_1\alpha_2\cdots\alpha_t$ t.q. $b(\alpha_j) = s(\alpha_{j+1})$, pour $j \in \{1,2,\cdots,t-1\}$.
		\end{itemize}
	\end{dfn}

	\begin{exem}
		Avec les carquois de l'exemple pr\'ec\'edant,
		\begin{enumerate}
			\item $\alpha, \alpha^2, \cdots$ sont des chemins dans $Q^1$,
			\item $\alpha, \beta, \gamma, \beta\alpha, \gamma\alpha$ sont des chemins dans $Q^2$.
		\end{enumerate}
	\end{exem}

	Soit $K$ un corps. $KQ$ est le $K$-espace vectoriel ayant pour base l'ensemble des chemins dans $Q$.

	\begin{exem}
		Avec les carquois de l'exemple pr\'ec\'edant,
		\begin{enumerate}
			\item La base de $KQ^1$ est $\{\varepsilon_1, \alpha, \alpha^2, \cdots\}$,

			$\varepsilon_1+3\alpha^2$ et $27\alpha^{28}-\alpha$ sont des \'el\'ements de $KQ^1$;
			\item La base de $KQ^2$ est $\{\varepsilon_1, \varepsilon_2, \varepsilon_3, \alpha, \beta, \gamma, \beta\alpha, \gamma\alpha\}$,

			$\alpha-\beta\alpha$ et $\varepsilon_1+\gamma+\alpha$ sont des \'el\'ements de $KQ^2$,

			$(\alpha-\beta\alpha) + (\varepsilon_1+\gamma+\alpha) = 2\alpha - \beta\alpha + \varepsilon_1 + \gamma$.
		\end{enumerate}
	\end{exem}

	La multiplication dans $KQ$ est d\'efinie comme suit

	Si $\alpha_1,\alpha_2,\cdots,\alpha_t, \beta_1,\beta_2,\cdots,\beta_s$ sont des chemins, alors on pose \[ (\alpha_1\alpha_2\cdots,\alpha_t) (\beta_1\beta_2\cdots\beta_s) = \left\{\begin{array}{lll}
		\alpha_1\cdots\alpha_t\beta_1\cdots\beta_s&\text{si}&b(\alpha_t) = s(\beta_1)\\
		0&\text{sinon}
	\end{array}\right\} \]

	\begin{exem}
		\begin{enumerate}
			\item $Q=Q^1$, $KQ = k[\alpha]$,
			\item $Q=Q^3$, la base de $KQ$ est $\{\varepsilon_1, \alpha, \alpha^2, \cdots, \beta, \beta^2, \cdots, \alpha\beta, \beta\alpha, \cdots\}$. Dans $KQ$, $\alpha\beta \neq \beta\alpha$, donc $KQ$ n'est pas $K[\alpha,\beta]$,
			\item $Q=Q^4$, la base de $KQ$ est $\{\varepsilon_1, \varepsilon_2, \varepsilon_3, \alpha, \beta, \beta\alpha\}$. $\varepsilon_1+\varepsilon_2+\varepsilon_3$ est neutre multiplicatif. En effet,
			\begin{align*}
				(\varepsilon_1+\varepsilon_2+\varepsilon_3)(\alpha+\beta)&= \varepsilon_1\alpha + \varepsilon_2\alpha + \varepsilon_3\alpha + \varepsilon_1\beta + \varepsilon_2\beta + \varepsilon_3\beta\\
				&= 0 + \alpha + 0 + 0 + 0 + \beta\\
				&= \alpha+\beta
			\end{align*}

			\[ \begin{bmatrix}
				K&0&0\\
				K&K&0\\
				K&K&K
			\end{bmatrix} = \left\{ \begin{bmatrix}
				e_1&0&0\\
				a&e_2&0\\
				c&b&e_3
			\end{bmatrix} ~\middle|~ e_1,e_2,e_3,a,b,c \in K \right\} \]

			On a une bijection \[ \begin{array}{rcl}
				\begin{bmatrix}
					K&0&0\\
					K&K&0\\
					K&K&K
				\end{bmatrix} &\to& KQ\\
				E_1=\begin{bmatrix}
					e_1&0&0\\
					0&0&0\\
					0&0&0
				\end{bmatrix} &\mapsto& \varepsilon_1\\
				&\vdots\\
				A=\begin{bmatrix}
					0&0&0\\
					1&0&0\\
					0&0&0
				\end{bmatrix} &\mapsto& \alpha\\
				B=\begin{bmatrix}
					0&0&0\\
					0&0&0\\
					0&1&0
				\end{bmatrix} &\mapsto& \beta\\
				C=\begin{bmatrix}
					0&0&0\\
					0&0&0\\
					1&0&0
				\end{bmatrix} &\mapsto& \beta\alpha
			\end{array} \] qui se comporte bien avec les produits
			\begin{align*}
				BA&= \begin{bmatrix}
					0&0&0\\
					0&0&0\\
					0&1&0
				\end{bmatrix} \begin{bmatrix}
					0&0&0\\
					1&0&0\\
					0&0&0
				\end{bmatrix}\\
				&= \begin{bmatrix}
					0&0&0\\
					0&0&0\\
					1&0&0
				\end{bmatrix}\\
				&= C
			\end{align*}
		\end{enumerate}
	\end{exem}

	\begin{dfn}
		Soit $A$ un anneau.
		\begin{enumerate}
			\item $a \in A_* = A \setminus \{0\}$ est un \emph{diviseur de $0$} s'il existe un autre \'el\'ement $b \in A_*$ t.q. $ab=0$,
			\item $A$ est dit \emph{int\`egre} s'il est commutatif et n'a pas de diviseur de $0$.
		\end{enumerate}
	\end{dfn}

	\begin{exem}
		\begin{enumerate}
			\item $\ZZ, \QQ, \RR, \ZZ[x], \QQ[x], \RR[x], \RR[x,y], \cdots$ sont tous int\`egres,
			\item $\ZZ_n$ est int\`egre si, et seulement si, $n$ est premier.

			En g\'en\'eral, les diviseurs de $0$ dans $\ZZ_n$ sont les $\overline{b}$ o\`u $\operatorname{pgcd}(b,n)>1$.
		\end{enumerate}
	\end{exem}

	\begin{rmq}
		Dans un anneau non int\`egre, la ``loi de simplification'' ne tient pas.
	\end{rmq}

	\begin{exem}
		Dans $\ZZ$,
		\begin{align*}
			x^2-3x+2&= 0\\
			(x-2)(x-1)&= 0\\
			x-2=0 \text{ ou } x-1=0\\
			x=2 \text{ ou } x=1
		\end{align*}

		Dans $\ZZ_6$,
		\begin{align*}
			x^2-\overline3x+\overline2&= 0\\
			(x-\overline2)(x-\overline1)&= 0
		\end{align*}
		On a des solutions: $x=\overline1, x=\overline2, x=\overline4, x=\overline5$.
	\end{exem}

	\begin{prop}
		Soit $A$ un anneau commutatif. Alors, $A$ est int\`egre si, et seulement si, la loi de simplification tient, c'est-\`a-dire, $ab=ac, a\neq0$ implique $b=c$.
		\begin{proof}
			\begin{subproof}[$\Rightarrow$:]
				Supposons que $A$ est int\`egre et soient $a,b,c \in A$ t.q. $a\neq0$ et $ab=ac$. On a alors $ab-ac=0$, donc $a(b-c)=0$. Comme $A$ est int\`egre et $a\neq0$, on a que $b-c=0$, donc $b=c$.
			\end{subproof}
			\begin{subproof}[$\Leftarrow$:]
				Supposons que $a\neq0$ et $a$ est un diviseur de $0$. Alors, $\exists b\neq0 \in A$ t.q. $ab=0$. Donc, $ab=a0$, alors $b=0$, une contradiction.
			\end{subproof}
		\end{proof}
	\end{prop}

	\begin{dfn}
		Soit $A$ un anneau. $a \in A$ est une \emph{unit\'e} s'il est inversible, c'est-\`a-dire, $\exists a'\in A$ t.q. $aa'=a'a=\id_A$. L'ensemble des unit\'es de $A$ est not\'e $\mathcal{U}(A)$.
	\end{dfn}

	\begin{exem}
		\begin{enumerate}
			\item $\mathcal{U}(\ZZ) = \{1,-1\}$,
			\item $\mathcal{U}(\QQ) = \QQ_*$,
			\item $\mathcal{U}(\RR[x]) = \RR_*$, les polyn\^omes de degr\'e $0$ non nuls,
			\item Soit $G$ un groupe et $A=K^{(G)}$, alors les \'el\'ements du groupe sont des unit\'es de $A$, $G \subseteq \mathcal{U}(A)$,
			\item $\mathcal{U}(\ZZ[i]) = ?$. $1,i,-1,-i \in \mathcal{U}(\ZZ[i])$. Supposons que $a+bi \in \mathcal{U}(\ZZ[i])$ et $c+di$ sont inverse. Alors, $(a+bi)(c+di)=1$, alors $\abs{a+bi} \abs{c+di}=1$, donc $(a^2+b^2)(c^2+d^2)=1$, d'o\`u $a^2+b^2=1$, donc $a^2=1$ ou $b^2=1$, ainsi $a=\pm1$ ou $b=\pm1$.
		\end{enumerate}
	\end{exem}

	\begin{dfn}
		Un corps est un anneau commutatif dont tout \'el\'ement non nul est une unit\'e.
	\end{dfn}

	\begin{exem}
		\begin{enumerate}
			\item $\QQ, \RR, \CC$ sont des corps;
			\item $\QQ[\sqrt{d}] = \left\{ x+y\sqrt{d} ~\middle|~ x,y \in \QQ \right\}$ est un anneau. Voyons que c'est un corps. Soit $x+y\sqrt{d} \in \QQ[\sqrt{d}]_*$. Dans $\CC$, on a $\dfrac{1}{x+y\sqrt{d}} = \dfrac{x-y\sqrt{d}}{x^2-dy^2} = \dfrac{x}{x^2-dy^2} + \left(\dfrac{-y}{x^2-dy^2}\right)\sqrt{d} \in \QQ[\sqrt{d}]$;
			\item $\ZZ_n$ est un corps si, et seulement si, $n$ est premier;
			\item Un corps non commutatif. Le corps des quaternions: $\mathbb{H} = \left\{ \begin{bsmallmatrix}
				a&b\\-\overline{b}&\overline{a}
			\end{bsmallmatrix} ~\middle|~ a,b\in\CC \right\}$. En effet,
			\begin{itemize}
				\item il s'agit d'un sous-anneau $\mathbb{H}\lte M_2(\CC)$
				\begin{itemize}
					\item Clairement un sous-groupe additif.
					\item $\begin{bsmallmatrix}
						1&0\\0&1
					\end{bsmallmatrix}\in\mathbb{H}$, car $\overline0=0$.
					\item Soient $\begin{bsmallmatrix}
						a&b\\-\overline{b}&\overline{a}
					\end{bsmallmatrix}, \begin{bsmallmatrix}
						c&d\\-\overline{d}&\overline{c}
					\end{bsmallmatrix}\in\mathbb{H}$. Alors
					\begin{align*}
						\begin{bmatrix}
							a&b\\-\overline{b}&\overline{a}
						\end{bmatrix} \begin{bmatrix}
							c&d\\-\overline{d}&\overline{c}
						\end{bmatrix}&= \begin{bmatrix}
							ac-b\overline{d}&ad+b\overline{c}\\
							-c\overline{b}-\overline{a}\overline{d}&-d\overline{b}+\overline{a}\overline{c}
						\end{bmatrix}\\
						&= \begin{bmatrix}
							ac-b\overline{d}&ad+b\overline{c}\\
							-\overline{(ad+b\overline{c})}&\overline{a}\overline{c}-b\overline{d}
						\end{bmatrix}
					\end{align*}

					De plus, $\det\begin{bsmallmatrix}
						a&b\\-\overline{b}&\overline{a}
					\end{bsmallmatrix} = a\overline{a}+b\overline{b} = \abs{a}^2+\abs{b}^2\neq0$, donc $\begin{bsmallmatrix}
						a&b\\-\overline{b}&\overline{a}
					\end{bsmallmatrix}$ est inversible dans $M_2(\CC)$. Or, \[ \begin{bmatrix}
					a&b\\-\overline{b}&\overline{a}
					\end{bmatrix}^{-1} = \frac{1}{\abs{a}^2+\abs{b}^2} \begin{bmatrix}
					\overline{a}&-b\\\overline{b}&a
					\end{bmatrix} \in \mathbb{H} \]
				\end{itemize}
			\end{itemize}
		\end{enumerate}
	\end{exem}

	\begin{rapp}[Inverses dans $\ZZ_p$]
		$p$ est premier.

		Soit $a \in \{0,1,\cdots,p-1\}$. Par le th\'eor\`eme de B\'ezout, $\exists s,t$ t.q. $as+bp=1$, donc $\overline{a} \overline{s} + \overline{b} \overline{p}=1$, c'est-\`a-dire $\overline{a} \overline{s}=1$. Ensuite, on utilise l'algorithme d'Euclide pour trouver $s$.
	\end{rapp}

	\begin{thm}
		Tout anneau int\`egre fini est un corps.
		\begin{proof}
			Soit $a \in A_*$ fix\'e. On veut montrer que $a$ est inversible.

			Posons $f_a:A \to A:t \mapsto at$. Si $t_1 \neq t_2$, on a $t_1-t_2\neq0$, d'o\`u $a(t_1-t_2)\neq0$, donc $f_a(t_1) \neq f_a(t_2)$, c'est-\`a-dire, $f_a$ est une injection.

			Comme $A$ est fini, $f_a$ est une surjection. Donc, $\exists x \in A$ t.q. $f_a(x)=\id$, c'est-\`a-dire, $ax=\id$.

			De m\^eme, $\exists x' \in A$ t.q. $x'a=\id$. $x' = x'\id = x'(ax) = (x'a)a = \id x = x$.
		\end{proof}
	\end{thm}

	\begin{rmq}
		On sait que $\forall n\in\NN$, il existe un anneau de cardinalit\'e $n$, notamment $\ZZ_n$ qui est un corps quand $n$ est premier.

		On peut se demander s'il existe des corps finis de cardinalit\'e arbitraire.

		$K$ un corps fini $\Longrightarrow \abs{K}=p^d$. De plus, $\forall p$ premier et $d\in\NN$, il existe un unique corps de cardinalit\'e $p^d$.
	\end{rmq}

	\begin{dfn}
		Soit $A$ un anneau. Sa \emph{caract\'eristique} est $\caract{A} = o(\id)$, l'ordre additif de $\id$.
	\end{dfn}

	\begin{exem}
		\begin{enumerate}
			\item $\caract{\ZZ} = 0 = \caract{\QQ} = \caract{\RR} = \caract{\CC}$;
			\item $\caract{\ZZ_n}=n$;
			\item $\caract{\ZZ_4\times\ZZ_6} = \operatorname{ppcm}(4,6) = 12$;
			\item en g\'en\'eral, $\caract{\ZZ_a\times\ZZ_b} = \operatorname{ppcm}(a,b)$.
			\begin{proof}[En effet]
				Soit $m=\operatorname{ppcm}(a,b)$, donc $m=aa'=bb'$. Notons $\overline{\id}$ l'identit\'e de $\ZZ_a$ et $\widetilde{\id}$ l'identit\'e de $\ZZ_b$,
				\begin{align*}
					m(\overline{\id}, \widetilde{\id})&= (m\overline{\id}, m\widetilde{\id})\\
					&= (aa'\overline{\id}, bb'\widetilde{\id})\\
					&= (a'\overline{a\id}, b'\widetilde{b\id})\\
					&= (\overline{0}, \widetilde{0})
				\end{align*}

				Si $m' \mid m$ est t.q. $m'(\overline{\id}, \widetilde{\id}) = (\overline{0}, \widetilde{0})$, c'est-\`a-dire $a \mid m'$ et $b \mid m'$ et donc $m \mid m'$.
			\end{proof}
		\end{enumerate}
	\end{exem}

	\begin{prop}
		Si $B \lte A$, alors $\caract{B} = \caract{A}$.
	\end{prop}

	\begin{thm}
		$A$ int\`egre, alors $\caract{A}=0$ ou $\caract{A}=p$ avec $p$ premier.
		\begin{proof}
			Soit $A$ int\`egre. Supposons $\caract{A}\neq0$.

			Si $\caract{A}=ab$, avec $a\neq1,b\neq1$, alors $0 = (ab)\id = (a\id) (b\id)$. L'int\'egrit\'e de $A$ donne une contradiction.
		\end{proof}
	\end{thm}

	\chapter{Id\'eaux et anneaux quotient}
	L'objectif est de transposer ``en anneaux'' la construction de groupe quotient, (homo)morphismes et isomorphismes.

	\begin{rapp}
		Soit $G$ un groupe et soit $H \lte G$ un sous-groupe.

		Dans $G$, on d\'efinit la congruence$\mod H$ par $g \equiv_H h \Leftrightarrow gh^{-1} \in H$ et $g \prescript{}{H}{\equiv} h \Leftrightarrow g^{-1}h \in H$.

		Les relations co\"incident si $H \trianglelefteq G$ est un sous-groupe normal de $G$, $ghg^{-1} \in H$, $\forall g \in G$, $\forall h \in H$.

		La classe de $g$ pour la congruence$\mod H$ est $gH = \left\{ gh ~\middle|~ h \in H \right\} = \overline{g}$.

		L'ensemble quotient $\sfrac{G}{\equiv_H}$ est not\'e $\sfrac{G}{H}$.

		Il s'agit d'un groupe pour $\overline{g_1} \; \overline{g_2} \coloneq \overline{g_1g_2}$.
	\end{rapp}

	Soit $(A,+,\cdot)$ un anneau. $(A,+)$ est un groupe ab\'elien additif. Alors, tout sous-groupe de $(A,+)$ est \emph{normal}.

	Si $I$ est un sous-groupe de $(A,+)$, on peut former sans probl\`eme le groupe quotient $\sfrac{A}{I}$ pour $(a+I)+(b+I) = (a+b)+I$.

	Pour ajouter une ``multiplication'', on pose $(a+I)(b+I) \coloneq (ab+I)$.

	Est-ce bien d\'efini? \[ \left. \begin{array}{rcl}
		a-a'&\in&I\\
		b-b'&\in&I
	\end{array}\right\} \overset{?}{\Rightarrow} ab-a'b' \in I \]

	\begin{exem}
		$(A,+,\cdot) = (\RR,+,\cdot)$, $(I,+,\cdot) = (\ZZ,+,\cdot)$. Le groupe quotient $\sfrac{\RR}{\ZZ}$ est connu. Voyons pour le produit: $a=\sqrt{2}, a'=1+\sqrt{2}, b=\sqrt{3}, b'=2+\sqrt{3}$. $a-a' \in \ZZ, b-b' \in \ZZ$.
		\begin{align*}
			ab-a'b'&= \sqrt{2}\sqrt{3}-(1+\sqrt{2})(2+\sqrt{3})\\
			&= -2-2\sqrt{2}-\sqrt{3}\\
			&\notin\ZZ
		\end{align*}

		Le probl\`eme est que le produit d'un \'el\'ement de $A$ par un \'el\'ement de $I$ n'appartient pas n\'ecessairement \`a $I$.
	\end{exem}

	\begin{lem}
		Soient $A$ un anneau et $I$ un sous-groupe additif de $A$. La congruence$\mod I$ est compatible avec le produit si, et seulement si, $x \in I, a \in A \Rightarrow ax,xa \in I$.
		\begin{proof}
			\begin{subproof}[$\Rightarrow$:]
				Supposons que $x \in I$, de sorte que $x-0 \in I$. Par ailleurs, pour $a \in A$, $a-a \in I$. D'o\`u $xa-0a \in I$, c'est-\`-dire $xa \in I$. De m\^eme, $ax \in I$.
			\end{subproof}
			\begin{subproof}[$\Leftarrow$:]
				Supposons que $a-a', b-b' \in I$. Alors,
				\begin{align*}
					a(b-b') + (a-a')b'&=  ab-ab'+ab'-a'b'\\
					&= ab-a'b'\\
					&\in I
				\end{align*}
			\end{subproof}
		\end{proof}
	\end{lem}

	\begin{dfn}[Id\'eal]
		Soient $A$ un anneau et $I \subseteq A$. On dit que $I$ est un \emph{id\'eal} de $A$, not\'e $I \trianglelefteq A$, si
		\begin{enumerate}
			\item $I$ est un sous-groupe additif de $A$,
			\item $x \in I, a \in A \Rightarrow ax,xa \in I$.
		\end{enumerate}
	\end{dfn}

	\begin{exem}
		\begin{enumerate}
			\item Soit $A$ un anneau, $0$ et $A$ sont des id\'eaux.
			\item Soit $I$ un id\'eal, $\id_A \in I \Rightarrow I=A$.

			En effet, pour $a \in A$, on a $a=\underbrace{a}_{\in A} \cdot \underbrace{\id_A}_{\in I} \in I$.

			De m\^eme, si $u \in \mathcal{U}(A)$ et $u \in I$, alors $I=A$.

			En effet, $\id_A=\underbrace{u}_{\in I} \cdot \underbrace{u^{-1}}_{\in A} \in I$.
			\item Quels sont les id\'eaux de $\ZZ$? Les sous-groupes de $\ZZ$ sont de la forme $n\ZZ$. Ce sont tous des id\'eaux.
			\item Quels sont les id\'eaux de $\ZZ_n$? $m\ZZ_n$.
			\begin{rapp}
				Les sous-groupes de $\ZZ_n = \sfrac{\ZZ}{n\ZZ}$ sont de la forme $\sfrac{G}{n\ZZ}$, o\`u $G \lte \ZZ$ t.q. $n\ZZ \subseteq G$.
			\end{rapp}

			$G = k\ZZ$, puis $n\ZZ \subseteq k\ZZ \Leftrightarrow k \mid n$.

			Les id\'eaux de $\ZZ_n$ sont donc de la forme $\sfrac{k\ZZ}{n\ZZ} = k\left(\sfrac{\ZZ}{n\ZZ}\right) = k\ZZ_n$ avec $k \mid n$.

			Concr\`etement, avec $n=6$. $A = \{\overline0, \overline1, \overline2, \overline3, \overline4, \overline5\}$. Si $k=1$, on a $A$. Si $k=2$, on a $2\ZZ_6 = \{\overline0, \overline2, \overline4\}$. Si $k=3$, on a $3\ZZ_6 = \{\overline0, \overline3\}$. Si $k=6$, on a $6\ZZ_6 = \{\overline0\}$.
			\item Soit $K$ un corps. Les seuls id\'eaux sont $0,K$.
			\item $A = T_2(\ZZ) = \left\{ \begin{bsmallmatrix}
				a&0\\b&c
			\end{bsmallmatrix} ~\middle|~ a,b,c\in\ZZ \right\}$ est un sous-anneau de $M_2(\ZZ)$. $I = \left\{ \begin{bsmallmatrix}
				0&0\\b&0
			\end{bsmallmatrix} ~\middle|~ b\in\ZZ \right\}$ est un sous-groupe additif de $T_2(\ZZ)$. Voyons que c'est un id\'eal.
			\begin{align*}
				\begin{bmatrix}
					a&0\\b&c
				\end{bmatrix} \begin{bmatrix}
					0&0\\x&0
				\end{bmatrix}&= \begin{bmatrix}
					0&0\\cx&0
				\end{bmatrix}\\
				&\in I\\
				\shortintertext{de m\^eme,}
				\begin{bmatrix}
					0&0\\x&0
				\end{bmatrix} \begin{bmatrix}
					a&0\\b&c
				\end{bmatrix}&\in I
			\end{align*}

			$I \trianglelefteq T_2(\ZZ)$.

			\begin{rmq}
				$I$ n'est pas un id\'eal de $M_2(\ZZ)$. En effet, $\begin{bsmallmatrix}
					a&b\\c&d
				\end{bsmallmatrix}\begin{bsmallmatrix}
					0&0\\x&0
				\end{bsmallmatrix} = \begin{bsmallmatrix}
					bx&*\\ *&*
				\end{bsmallmatrix} \notin I$.
			\end{rmq}
			\item $A=M_2(\QQ)$ n'a aucun id\'eal propre. En effet, supposons $M = \begin{bsmallmatrix}
				a&b\\c&d
			\end{bsmallmatrix} \in I \trianglelefteq A$, $M\neq0$. On peut supposer, sans perte de g\'en\'eralit\'e, que $a\neq0$.En effectuant des op\'erations \'el\'ementaires de lignes et de colonnes, on obtient $\begin{bsmallmatrix}
				1&0\\0&*
			\end{bsmallmatrix} \in I$. Si $\begin{bsmallmatrix}
				1&0\\0&1
			\end{bsmallmatrix} \in I$, alors $I=A$. Si $\begin{bsmallmatrix}
				1&0\\0&0
			\end{bsmallmatrix} \in I$, alors $\begin{bsmallmatrix}
				0&0\\0&1
			\end{bsmallmatrix} \in I$, et donc $\begin{bsmallmatrix}
				1&0\\0&1
			\end{bsmallmatrix} \in I$.

			En g\'en\'eral, $M_n(K)$ n'a aucun id\'eal propre.
		\end{enumerate}
	\end{exem}

	\begin{thm}
		Soient $A$ un anneau, $(I_i)_{i \in I}$ une famille arbitraire d'id\'eaux de $A$. Alors, $\bigcap\limits_{i \in I} I_i$ est un id\'eal de $A$.
		\begin{proof}
			Exercice
		\end{proof}
	\end{thm}

	\begin{prop}[et d\'efinition]
		Soient $A$ un anneau, $X \subseteq A$ une partie de $A$ et $\mathscr{F}_X = \{I \trianglelefteq A \mid X \subseteq I\}$.

		Alors, l'id\'eal $\bigcap\limits_{I \in \mathscr{F}_X}I$, not\'e $\langle X \rangle$, v\'erifie
		\begin{itemize}
			\item $X \subseteq \langle X \rangle$,
			\item si $X \subseteq J$ et $J \trianglelefteq A$, alors $\langle X \rangle \subseteq J$.
		\end{itemize}
		C'est l'\emph{id\'eal de $A$ engendr\'e par $X$}.
		\begin{proof}
			$\mathscr{F}_X\neq\emptyset$, car $A \in \mathscr{F}_X$. Par le th\'eor\`eme pr\'ec\'edent, $\bigcap\limits_{I \in \mathscr{F}_X}I \trianglelefteq A$.

			Pour $I \in \mathscr{F}_X$, on a $X \subseteq I$, d'o\`u $X \subseteq \bigcap\limits_{I \in \mathscr{F}_X}I$.

			Si $X \subseteq J$, $J \trianglelefteq A$, on a $J \in \mathscr{F}_X$, donc $\bigcap\limits_{I \in \mathscr{F}_X}I \subseteq J$.
		\end{proof}
	\end{prop}

	\begin{exem}
		\begin{enumerate}
			\item $\langle \emptyset \rangle = 0 = \langle 0 \rangle$.
			\item $\langle \id_A \rangle = A$.
			\item Dans $A=\ZZ$, $X=\{a,b\}$. $\langle a,b \rangle = d\ZZ$, o\`u $d=\operatorname{pgcd}(a,b)$.

			En effet, $\langle a,b \rangle$ est un id\'eal de $\ZZ$, donc de la forme $k\ZZ$, pour un $k\in\ZZ$. $\langle a,b \rangle \supseteq \langle a \rangle = a\ZZ$. $\langle a,b \rangle \supseteq \langle b \rangle = b\ZZ$. Donc, $a\ZZ \subseteq k\ZZ$ et $b\ZZ \subseteq k\ZZ$. Alors, $k \mid a$ et $k \mid b$.
			\begin{rapp}
				$d$ est un $\operatorname{pgcd}$ de $a,b$ si
				\begin{enumerate}
					\item $d \mid a$ et $d \mid b$,
					\item si $d' \mid a$ et $d' \mid b$, alors $d' \mid d$.
				\end{enumerate}
			\end{rapp}

			De plus, si $k' \mid a$ et $k' \mid b$, alors $a=k$ et $b=k'b'$ et donc on a $a,b \in k'\ZZ$, d'o\`u $k\ZZ \subseteq k'\ZZ$. Ainsi, $k' \mid k$ et $k$ est $\operatorname{pgcd}(a,b)$.
		\end{enumerate}
	\end{exem}

	Si $I_1,I_2 \trianglelefteq A$, que peut-on dire de $I_1 \cup I_2$?
	\begin{exem}
		$A=\ZZ, I_1=3\ZZ, I_2=4\ZZ$. Alors, $1=4-3 \in \langle 3\ZZ \cup 4\ZZ \rangle$, mais $1 \notin 3\ZZ \cup 4\ZZ$.

		En g\'en\'eral, $I_1 \cup I_2$ n'est pas un id\'eal de $A$.
	\end{exem}

	\begin{dfn}
		Soient $A$ un anneau et ,$I,J \trianglelefteq A$ des id\'eaux. On d\'efinit $I+J \coloneq \left\{ x+y ~\middle| x\in I, y\in J \right\}$.
	\end{dfn}

	\begin{lem}
		Soient $A$ un anneau et $I,J \trianglelefteq A$ des id\'eaux.

		\begin{enumerate}
			\item $I+J \trianglelefteq A$,
			\item $I+J = \langle I \cup J \rangle$.
		\end{enumerate}
		\begin{proof}
			\begin{enumerate}
				\item $I+J \neq \emptyset$, car $0=\underbrace{0}_{\in I}+\underbrace{0}_{\in J}$.

				Supposons que $a=x+y$, $b=x'+y'$ avec $x,x'\in I$ et $y,y'\in J$. Alors
				\begin{align*}
					a-b&= (x+y)-(x'+y')\\
					&= \underbrace{(x-x')}_{\in I}+\underbrace{(y-y')}_{\in J}\\
					&\in I+J
				\end{align*}

				Soient $a=x+y\in I+J$ et $\alpha\in A$.
				\begin{align*}
					\begin{split}
						\alpha a&= \alpha(x+y)\\
						&= \underbrace{\alpha x}_{\in I}+\underbrace{\alpha y}_{\in J}\\
						&\in I+J
					\end{split}
					&
					\begin{split}
						a\alpha&= (x+y)\alpha\\
						&= x\alpha+y\alpha\\
						&\in I+J
					\end{split}
				\end{align*}
				\item Si $x\in I$, alors $x=x+\underbrace{0}_{\in J}\in I+J$ de sorte que $I\subseteq I+J$. De m\^eme, $J\subseteq I+J$. D'o\`u, $I \cup J \subseteq I+J$.

				Supposons que $K \trianglelefteq A$ t.q. $I\cup J\subseteq K$. Alors, un \'el\'ement $z=x+y\in I+J$ v\'erifie $z\in K$ puisque $x\in I\subseteq K$ et $y\in J\subseteq K$.

				Ainsi, $I+J\subseteq K$.
			\end{enumerate}
		\end{proof}
	\end{lem}

	On peut se demander comment d\'ecrire explicitement $\langle X \rangle$ en g\'en\'eral, avec $X \subseteq A$.

	\begin{lem}
		Soient $A$ un anneau et $X\subseteq A$ une partie. Alors, \[\langle X \rangle = AXA = \left\{ \sum_{i=1}^{n}a_ix_ib_i ~\middle|~ a_i,b_i\in A, x_i\in X \right\} \]
		\begin{proof}
			\begin{itemize}
				\item $AXA\neq\emptyset$.

				Supposons que $\sum_i a_ix_ib_i, \sum_i c_iy_id_i \in AXA$. Alors, $\sum_i a_ix_ib_i - \sum_i c_iy_id_i \in AXA$.

				Supposons que $\sum_i a_ix_ib_i \in AXA$, $\alpha \in A$. Alors,
				\begin{align*}
					\alpha\left(\sum_i a_ix_ib_i\right)&= \sum_i (\alpha a_i)x_ib_i\\
					&\in AXA
				\end{align*}
				De m\^eme, $\left(\sum_ia_ix_ib_i\right)\alpha\in AXA$.

				Ainsi, $AXA\trianglelefteq A$.

				De plus, si $x\in X$, $x=\id x\id\in AXA$, d'o\`u $X\subseteq AXA$.
				\item Supposons que $J\trianglelefteq A$ t.q. $X\subseteq J$. Tout \'el\'ement de la forme $a_ix_ib_i$ avec $a_i,b_i\in A$, $x_i\in X$ se trouve dans $J$, car $J$ est un id\'eal. Comme $J$ est aussi un sous-groupe additif de $A$, on a que $AXA \subseteq J$.
			\end{itemize}
		\end{proof}
	\end{lem}

	\begin{rmq}
		Si $A$ est commutatif, alors $\langle X \rangle = AX = XA$. Il s'agit de l'ensemble des combinaisons lin\'eaires d'\'el\'ements de $X$ \`a coefficients dans $A$.
	\end{rmq}

	\begin{exer}
		Soit $A$ un anneau commutatif. $a \in A$.
		\begin{proof}[m.q.]
			$\langle a \rangle = A \Longleftrightarrow a\in \mathcal{U}(A)$.
			\begin{subproof}[$\Rightarrow$:]
				Si $\langle a \rangle = A$, alors $\id\in\langle a\rangle=Aa$. D'o\`u l'existence de $u\in A$ t.q. $\id=ua$, c'est-\`a-dire, $a\in\mathcal{U}(A)$.
			\end{subproof}
			\begin{subproof}[$\Leftarrow$:]
				Si $a\in\mathcal{U}(A)$, alors $a^{-1}\in A$ et donc $\id=aa^{-1}\in\langle a\rangle$. D'o\`u $\langle a\rangle=A$.
			\end{subproof}
		\end{proof}
	\end{exer}

	\begin{rmq}
		Dans $\ZZ$, $\langle3,11\rangle = \langle\operatorname{pgcd}(3,11)\rangle = \langle1\rangle = \ZZ$.

		$\langle3,11\rangle$ est engendr\'e par un unique \'el\'ement.
	\end{rmq}

	\begin{dfn}
		Soit $A$ un anneau commutatif. Un id\'eal $I\trianglelefteq A$ est dit \emph{principal} s'il est de la forme $I=Ax=\langle x\rangle$. On dit que $x$ est \emph{g\'en\'erateur}.
	\end{dfn}

	\begin{prop}
		Soit $A$ un anneau commutatif. Alors, $A$ est un corps si, et seulement si, ses seuls id\'eaux sont $0$ et $A$.
		\begin{proof}
			\begin{subproof}[$\Rightarrow$:]
				D\'ej\`a fait.
			\end{subproof}
			\begin{subproof}[$\Leftarrow$:]
				Si $x\neq0\in A$, alors $\langle x\rangle=A$, d'o\`u $x$ est inversible.
			\end{subproof}
		\end{proof}
	\end{prop}
\end{document}
